% ********************************************
\section{Long-Term Waiting List Management\label{sec:waitinglist}}
% *******************************************

As discussed in Section~\ref{sec:lit_review}, CMP constitutes the strategic decision layer responsible for determining the volumes of elective patients to be admitted over a given planning horizon.
In contrast to much of the existing literature, which primarily selects these volumes to optimise financial or capacity-related objectives, this work adopts a stability-oriented perspective.
We aim to enable decision-makers to impose maximum acceptable waiting times (e.g., six months) and to guarantee, in the long run, that such targets are met without allowing waiting lists to grow unboundedly. 
This is particularly relevant in the context of public healthcare systems, such as Brazil's Unified Health System (SUS), where there is a legal obligation to provide care to all patients in need.
Thus, ensuring equitable access to care and adherence to waiting time standards is paramount.

In line with this objective, this section introduces the theoretical foundations of the modified $(R,Q)$ control policy proposed to manage elective surgery waiting lists under stochastic demand.
We first characterise the dynamics of the underlying stochastic process and analyse its long-run behaviour. 
Subsequently, we derive analytical upper and lower bounds on the batch size $Q$ that ensures compliance with prescribed waiting-time requirements, thereby establishing the strategic input for the integrated planning framework.
Table~\ref{tab:symb_rq} summarises the notation used in this section.

\begin{table}[htb]
    \centering
    \caption{Symbols used in the long-term waiting list management model.\label{tab:symb_rq}}
    \scriptsize
    \begin{tabular}{cl}
    \toprule
    \textbf{Index} & \textbf{Description} \\
    \midrule
    $\stocprocperiod$ & Planning period index \\
    $\stocprocremainstate$ & State index of the remainder process \\
    $i,j$ & Generic state indices \\
    \midrule

    \textbf{Set} & \textbf{Description} \\
    \midrule
    $\statespace$ & State space of process $\stocprocqueue{\stocprocperiod}$ \\
    $\statespace_{\stocprocbatches{}}$ & State space of process $\stocprocbatches{\stocprocperiod}$ \\
    $\statespace_{\stocprocremain{}}$ & State space of process $\stocprocremain{\stocprocperiod}$ \\
    $\Omega$ & Support of the demand random variable \\
    \midrule

    \textbf{Parameter} & \textbf{Description} \\
    \midrule
    $\rthreshold$ & Reorder threshold of the modified $(R,Q)$ policy \\
    $\batchsizeq$ & Batch size under the modified $(R,Q)$ policy \\
    $\waittimereq$ & Waiting-time target \\
    $\conflevel$ & Probability threshold for the waiting-time requirement \\
    $\upq$ & Upper bound on feasible batch sizes \\
    $\lowq$ & Lower bound on feasible batch sizes \\
    $\maxnumberbatch$ & Maximum number of batch arrivals in one period \\
    \midrule

    \textbf{Random Variable / Stochastic Process} & \textbf{Description} \\
    \midrule
    $\stocprocqueue{\stocprocperiod}$ & Queue-length process \\
    $\rvdemand$ & Demand between two consecutive planning periods \\
    $\drawnrvdemand$ & Realisation of $\rvdemand$ \\
    $\stocprocbatches{\stocprocperiod}$ & Batch-level queue process \\
    $\stocprocremain{\stocprocperiod}$ & Remainder process associated with $\stocprocqueue{\stocprocperiod}$ \\
    $\roundbatches$ & Number of batch arrivals in one period \\
    $\omega$ & Waiting time of an arriving patient \\
    \midrule

    \textbf{Function} & \textbf{Description} \\
    \midrule
    $f_{\rvdemand}(\cdot)$ & Probability mass function of $\rvdemand$ \\
    $F_{\rvdemand}(\cdot)$ & Cumulative distribution function of $\rvdemand$ \\
    $\roundbatches(\stocprocremainstate)$ & Number of batch arrivals given remainder state $\stocprocremainstate$ \\
    $\steadystate{\stocprocremain{}}(\cdot)$ & Steady-state distribution of $\stocprocremain{\stocprocperiod}$ \\
    $\steadystate{\stocprocbatches{}}(\cdot)$ & Steady-state distribution of $\stocprocbatches{\stocprocperiod}$ \\
    $\transmatrix{\stocprocremain{}}$ & Transition matrix of process $\stocprocremain{\stocprocperiod}$ \\
    $\transmatrix{}^{\stocprocbatches{}}$ & Transition matrix of process $\stocprocbatches{\stocprocperiod}$ \\
    $\transmatrixval{\stocprocremainstate\stocprocremainstate'}{\stocprocremain{}}(k)$ & Conditional transition probability of $\stocprocremain{\stocprocperiod}$ given demand $k$ \\
    $\transmatrixval{ij}{}$ & Generic transition probability entry \\
    \midrule

    \textbf{Derived Quantity} & \textbf{Description} \\
    \midrule
    $\1_{\{\stocprocqueue{\stocprocperiod} \ge \rthreshold\}}$ & Indicator of batch processing \\
    \bottomrule
    \end{tabular}
\end{table}

% *******************************************
\subsection{$(R,Q)$ surgery scheduling policy: long-term distributions\label{sec:rqtheory}}
% *******************************************

Let $\stocprocqueue{\stocprocperiod}, \, \stocprocperiod \ge 0$ be a stochastic process representing an elective surgery queue with stationary and independent increments over discrete planning periods, with $\stocprocqueue{\stocprocperiod}$ representing the number of patients in the system at period $\stocprocperiod \ge 0$. Assume that the queue is controlled via a modified $(R,Q)$ policy \citep[e.g.,][]{Axsater2015Inv} whereby a batch of $\batchsizeq \in \mathbb N$ surgeries is processed at time period $\stocprocperiod$ whenever $\stocprocqueue{\stocprocperiod} \ge \rthreshold, \, \rthreshold \ge \batchsizeq$ at the start of the period, where $\mathbb N$ is the set of positive integers. In contrast to the classical $(R,Q)$ policy, which processes as many batches as necessary to bring the inventory below the reorder point ($R$), our policy can only process a batch at a time and relies on long-term properties to ensure that the process $\stocprocqueue{\stocprocperiod}, \, \stocprocperiod \ge 0$ is stable and guarantees prescribed waiting time requirements.

Let $\rvdemand$ be a non-negative integer random variable representing the overall demand between two consecutive planning periods, with probability distribution $f_{\rvdemand}: \Omega \to [0,1), \, \Omega \subset \mathbb Z_+$, and cumulative distribution $F_{\rvdemand}: \mathbb Z_+ \to [0,1]$ given by
%
\[
F_{\rvdemand}(d) = \sum_{l=0}^{\drawnrvdemand} f_{\rvdemand}(l),
\]
%
where $\mathbb Z_+$ denotes the set of non-negative integers. This implies that the dynamics of process $\stocprocqueue{\stocprocperiod}, \, \stocprocperiod \ge 0$ can be expressed as:
%
\begin{equation} \label{eq:dynamics}
\stocprocqueue{\stocprocperiod+1} = \stocprocqueue{\stocprocperiod} + \drawnrvdemand - \batchsizeq \1_{\{\stocprocqueue{\stocprocperiod} \ge \rthreshold\}},
\end{equation}
%
where $\1_{\{A\}}$ is the indicator function, which is equal to $1$ if $A$ holds and is nil otherwise, and $\drawnrvdemand$ is a value drawn from $f_{\rvdemand}$ that represents the surgical demand in period $\stocprocperiod$. The state space of process $\stocprocqueue{\stocprocperiod}, \, \stocprocperiod \ge 0$ is given by $\statespace = \{ 0, \, 1 , \, 2, \, \ldots \}$, where $\stocprocqueue{\stocprocperiod}=i$ indicates that there are $\rthreshold-\batchsizeq + i$ patients waiting for elective surgery.

Letting 
%
\begin{equation} \label{eq:defYZ}
\stocprocbatches{\stocprocperiod} = \left \lfloor \dfrac{\stocprocqueue{t}}{Q} \right \rfloor, \qquad Z_t = \stocprocqueue{t} \, \% \,  \batchsizeq,
\end{equation}
%
where $\%$ represents the modulus operator, we can write:
%
\[
\stocprocqueue{\stocprocperiod} = Q \stocprocbatches{\stocprocperiod} + \stocprocremain{\stocprocperiod},
\]
% 
where $\stocprocbatches{\stocprocperiod}, \, \stocprocperiod \ge 0$ is the integer number of batches required to bring the queue instantaneously below $\rthreshold$ and $\stocprocremain{\stocprocperiod} \in \{0, 1, \ldots, \batchsizeq-1\}$ is the number of patients that would remain in excess of the minimum allowed queue of $(\rthreshold-\batchsizeq)$ patients after removing $\stocprocbatches{\stocprocperiod}$ batches of $\batchsizeq$ patients instantaneously.

Note that, from the definitions above, the state space of process $\stocprocbatches{\stocprocperiod}, \, \stocprocperiod \ge 0$ is $\statespace_{\stocprocbatches{}} = \mathbb Z_+$, while $\statespace_{\stocprocremain{}} = \{0, \, 1, \, \ldots, \, Q-1\}$ is the state space of process $\stocprocremain{\stocprocperiod}, \, \stocprocperiod \ge 0$. In the next result, we will derive the long-term distribution of process $\stocprocremain{\stocprocperiod}, \, \stocprocperiod \ge 0$.

\begin{prop} \label{prop:steadystate}
Let $\steadystate{\stocprocremain{}}: \statespace_{\stocprocremain{}} \to [0,1]$ be the steady state distribution of process $\stocprocqueue{\stocprocperiod}, \, \stocprocperiod \ge 0$. Then, it holds that 
%
\begin{equation} \label{eq:steadyZ}
\steadystate{\stocprocremain{}}(\stocprocremainstate) = \dfrac{1}{\batchsizeq}, \forall \stocprocremainstate \in \statespace_{\stocprocremain{}}.
\end{equation}
%
\end{prop}
%
\begin{proof}
The proof closely follows that of \citet[Proposition 5.1]{Axsater2015Inv}. Let $\transmatrix{\stocprocremain{}}$ be the transition matrix describing the dynamics of process $\stocprocremain{\stocprocperiod}, \, \stocprocperiod \ge 0$. Assume that $\stocprocremain{\stocprocperiod} = \stocprocremainstate, \, \stocprocremainstate \in \statespace_{\stocprocremain{}}$ at the start of period $\stocprocperiod \ge 0$, and let and $\rvdemand=k$ be the number of new patient arrivals in period $\stocprocperiod$. Hence, it follows that 
%
\[
\stocprocremain{\stocprocperiod+1} =  (\stocprocremainstate + k) \,\%\, \batchsizeq  = (\stocprocremainstate + k + \batchsizeq \stocprocbatches{\stocprocperiod}) \,\%\, \batchsizeq  = (\stocprocqueue{\stocprocperiod} + k) \,\%\, \batchsizeq,
\] 
%
where the last equality come from the definitions in Eq.~\eqref{eq:defYZ}. Therefore, given $\rvdemand=k$, it is evident that $\transmatrixval{\stocprocremainstate\stocprocremainstate'}{\stocprocremain{}}(k) := P( \stocprocremain{\stocprocperiod+1} = \stocprocremainstate' | \stocprocremain{\stocprocperiod} = \stocprocremainstate, \rvdemand=k )$ equals one for exactly one value of $\stocprocremainstate$ and zero otherwise. Consequently, we can write:
%
\[
\sum_{\stocprocremainstate=0}^{\batchsizeq-1} \transmatrixval{\stocprocremainstate\stocprocremainstate'}{\stocprocremain{}} = \sum_{\stocprocremainstate=0}^{\batchsizeq-1} \mathbb{E}_{k} \left\{ \transmatrixval{\stocprocremainstate\stocprocremainstate'}{\stocprocremain{}}(k) \right\} = \mathbb{E}_k \left\{  \sum_{\stocprocremainstate=0}^{\batchsizeq-1} \transmatrixval{\stocprocremainstate\stocprocremainstate'}{\stocprocremain{}}(k) \right\} = \mathbb{E}_k \{ 1 \} = 1.
\]
%
The expression above  implies that the Markov chain describing the dynamics of process $Z_t, \, t \ge 0$ is doubly stochastic, and therefore has a uniform steady state distribution. Indeed, substituting \eqref{eq:steadyZ} in the steady state equations:
%
\[
\steadystate{}^{\stocprocremain{}}(\stocprocremainstate) = \sum_{\stocprocremainstate'=0}^{\batchsizeq-1} \steadystate{}^{\stocprocremain{}}(\stocprocremainstate') \transmatrixval{\stocprocremainstate'\stocprocremainstate}{\stocprocremain{}},
\]
%
one can easily see that the steady state distribution satisfies Eq. \eqref{eq:steadyZ}.
%
\end{proof}

The result in Proposition \ref{prop:steadystate} is very important, as it shows that the steady state distribution of process $\stocprocremain{\stocprocperiod}, \, \stocprocperiod \ge 0$, is uniform regardless of the distribution of patient arrivals. Next, we will use this fact to derive the steady state distribution of process $\stocprocbatches{\stocprocperiod}, \, \stocprocperiod \ge 0$. 

% ********************************************
\subsection{The dynamics of process $\stocprocbatches{\stocprocperiod}, \, \stocprocperiod \ge 0$\label{sec:batchprocess}}
% *******************************************

Assume that process $\stocprocremain{\stocprocperiod}, \, \stocprocperiod \ge 0$ has reached equilibrium, which is also equivalent to supposing that $P(\stocprocremain{0} = \stocprocremainstate) = \steadystate{\stocprocremain{}}(\stocprocremainstate) = \dfrac{1}{\batchsizeq}$ \citep[e.g.,][]{Bremaud2020}. We can describe the dynamics of process $\stocprocbatches{\stocprocperiod}, \, \stocprocperiod \ge 0$ as:
%
\begin{equation} \label{eq:dynY}
\stocprocbatches{\stocprocperiod+1} = \begin{cases}
\stocprocbatches{\stocprocperiod} - 1 + \roundbatches, & \text{if} \; \stocprocbatches{\stocprocperiod} \ge 0, \\
\roundbatches; \; & \text{if} \; \stocprocbatches{\stocprocperiod} = 0,
\end{cases}
\end{equation}
%
where $\roundbatches$ is the round number of batch arrivals at time $\stocprocperiod$. To fully characterize the random variable $\roundbatches$, it is necessary to define
%
\begin{equation} \label{eq:Az}
\roundbatches(\stocprocremainstate) = \left \lfloor \dfrac{\rvdemand + \stocprocremainstate}{\batchsizeq} \right \rfloor,
\end{equation}
%
which represents the round number of batch arrivals at time $\stocprocperiod$, given that $\stocprocremain{\stocprocperiod} = \stocprocremainstate$. Then, in view of Proposition~\ref{prop:steadystate}, it follows that:
%
\begin{equation} \label{eq:batcharr}
P(\roundbatches = k ) = \dfrac{1}{\batchsizeq} \sum_{\stocprocremainstate=0}^{\batchsizeq-1} P(\roundbatches(\stocprocremainstate) = k ).    
\end{equation}

Therefore, from \eqref{eq:dynY} we have:
%
\begin{equation} \label{eq:ydynamicY}
\transmatrixval{ij}{\stocprocbatches{}} = P(\stocprocbatches{\stocprocperiod+1} = j | \stocprocbatches{\stocprocperiod} = i ) = \begin{cases}
    P\left( \roundbatches = j  \right), & \text{if} \; \stocprocbatches{\stocprocperiod} = 0, \\
    P\left( \roundbatches = j-i + 1 \right),  & \text{if} \; \stocprocbatches{\stocprocperiod} \ge 1.
\end{cases}
\end{equation}
%

Note that Eq. \eqref{eq:dynY}-\eqref{eq:ydynamicY} correspond to the dynamics in \citet[Section 6.1.2]{Shortle2018}. Similarly to \citet{Shortle2018}, the transition matrix of process $\stocprocbatches{\stocprocperiod}, \, \stocprocperiod \ge 0$ can be written as:
%
\begin{equation}
\transmatrix{}^{\stocprocbatches{}} = [\transmatrixval{ij}{}] = \left[
\begin{array}{cccccccc}
k_0 & k_1 & k_2 & \ldots & k_M & 0 & 0 & \ldots \\
k_0 & k_1 & k_2 & \ldots & k_M & 0 & 0 & \ldots \\
0 & k_0 & k_1 & k_2 & \ldots & k_M & 0 & \ldots \\
0 & 0 & k_0 & k_1 & k_2 & \ldots & k_M & \ldots \\
\vdots & \vdots & \vdots & \vdots & \vdots & \vdots & \vdots & 
\end{array}
\right],
\end{equation}
%
where $\maxnumberbatch$ is the maximum number of new batch arrivals, i.e., the maximum value that random variable $\roundbatches$ can assume.

The steady state probabilities of process $\stocprocbatches{\stocprocperiod}, \, \stocprocperiod \ge 0$ can then be recursively evaluated by solving:
\begin{equation} \label{eq:ssY}
\steadystate{\stocprocbatches{}}^T = \steadystate{\stocprocbatches{}}^T P,
\end{equation}
%
with the contour condition that $\steadystate{\stocprocbatches{}}(0) = 1 - \rho $ \citep[Eq. (6.12)]{Shortle2018}, 
%
%\begin{equation} \label{eq:ssY}
%\pi_Y(i) = \begin{cases}
%1 - \rho, & \text{if} \; i = 0, \\
%(1 - \rho) k_i + \displaystyle \sum_{j=1}^{i+1} \pi_Y(j) k_{i-j+1}, \, i \ge 0
%\end{cases}
%\end{equation}
%
where $\rho = E(\roundbatches)$. Observe that, to guarantee stability, we must have $E(\roundbatches) < 1$. Note also that the steady state probabilities up to a given state $j$ can be recursively calculated using the first $j$ equations in \eqref{eq:ssY}.

% ********************************************
\subsection{Upper and lower bounds on the batch size $\batchsizeq$ under waiting time requirements} \label{sec:bounds}
% *******************************************

Let $\omega$ be the waiting time of a patient joining the queue, and assume that the healthcare system requires that:
%
\begin{equation} \label{eq:waitthreshold}
P( \omega > \tau ) < \alpha,
\end{equation}
%
for some target waiting time $\waittimereq >0$ and probability threshold $\conflevel \in (0,1)$. Now, if a patient arrives and encounters $\stocprocbatches{\stocprocperiod} > 0$ batches of patients in the queue, their turn will only arrive once these $\stocprocbatches{\stocprocperiod}$ batches are processed, which will take $\stocprocbatches{\stocprocperiod}$ time periods. This means that the patient will have to wait $\waittimereq = \stocprocbatches{\stocprocperiod}$ periods for their turn. Therefore, Eq. \eqref{eq:waitthreshold} can be equivalently stated as:
%
\begin{equation} \label{eq:waitthresholdy}
\sum_{i = 0}^\waittimereq \steadystate{\stocprocbatches{}}(i) > 1 - \conflevel,
\end{equation}
%
where $\steadystate{\stocprocbatches{}} : \statespace_{\stocprocbatches{}} \to (0,1)$ satisfies Eq. \eqref{eq:ssY}. Therefore, to limit the waiting time according to the prescribed requirements, we can establish a lower bound ($\lowq$) on the batch size $\batchsizeq > 0$ as follows:
%
\begin{equation} \label{eq:under}
\lowq = \min \{ \batchsizeq \in \mathbb N : \sum_{i = 0}^\waittimereq \steadystate{\stocprocbatches{}}(i) > 1- \conflevel \}.
\end{equation}

When $\stocprocbatches{\stocprocperiod} = 0$, on the other hand, the next batch will only be processed when the system leaves this state, therefore by applying the strong Markov property \citep{Bremaud2020}, we have:
%
\[
\omega = \min \{\delta > 0 : \stocprocbatches{\delta} > 0 | \stocprocbatches{0} = 0 \}.
\]
%
Considering the transition matrix in Eq.~\eqref{eq:ydynamicY}, it follows that $\omega$ will be a geometrically distributed variable and:
%
\begin{equation} \label{eq:upperbound}
P( \omega > \waittimereq | \stocprocbatches{0} = 0 )  = k_0^{\waittimereq}.
\end{equation}

One can intuitively see that the value of $k_0$ will increase as $\batchsizeq$ increases, since in this case we will need more patient arrivals until we are able to process a full batch of patients. Hence, from Eq. \eqref{eq:upperbound}, we can set up an upper bound $\upq$ for the batch size such that:
%
\begin{equation} \label{eq:overbound}
\upq = \max \{ \batchsizeq \in \mathbb N : k_0^{\waittimereq} < \conflevel \}.
\end{equation}


% Therefore, considering a set of medical subspecialties $\specialtyset$, we can calculate batch sizes $\upq_{\specialtyidx}$ and $\lowq_{\specialtyidx}$ for every specialty $\specialtyidx \in \specialtyset$, and then ensure that the waiting time requirements are satisfied both when patients arrive to an empty queue and when they encounter batches of patients in the queue.
In the example below, we show how to derive general bounds $\lowq$ and $\upq$ for a given demand distribution and waiting time requirements.

\begin{ex}[Deriving lower and upper bounds for $\batchsizeq$]
Consider an elective surgery queue with monthly demand $\rvdemand$ that satisfies the distribution in Table \ref{tab:demprob}.
%
\begin{center}
\begin{table}[h!]
\caption{Monthly demand distribution for surgeries} \label{tab:demprob}
\scriptsize{
\begin{tabular}{|l|c|c|c|c|c|c|c|c|c|c|c|} \hline \hline
\textbf{Demand ($\drawnrvdemand$)} & 0 & 1 & 2 & 3 & 4 & 5 & 6 & 7 & 8 & 9 & 10 \\ \hline
$P(\rvdemand=\drawnrvdemand)$ & 0.0308 & 0.0393 & 0.1179 & 0.2096 & 0.2445 & 0.1956 & 0.1086 & 0.0414 & 0.0103 & 0.0015 & 0.0005 \\ \hline
\end{tabular}
}
\end{table}
\end{center}
%
The expected monthly demand is therefore $\mathbb{E}(\rvdemand) = 3.9022$. Setting $\batchsizeq = 4$, which ensures long-term stability, we have $\statespace_{\stocprocremain{}{}} = \{0, 1, 2, 3\}$. This allows us to evaluate the individual distributions of $\roundbatches(\stocprocremainstate), \, \stocprocremainstate \in \statespace_{\stocprocremain{}}$ as defined in Eq. \eqref{eq:Az}, which we depict in the Table \ref{tab:dembatches}, along with the distribution of variable $\roundbatches$ from Eq. \eqref{eq:batcharr}.


\begin{table}[h!]
\centering \caption{Monthly demand distribution for surgeries} \label{tab:dembatches}
\scriptsize{
\begin{tabular}{|l|c|c|c|c|} \hline \hline
$P(\roundbatches(\stocprocremainstate) = \drawnrvdemand) $ & \multicolumn{4}{|c|}{\textbf{Batch demand ($\drawnrvdemand$)}} \\ \cline{2-5}
  & 0 & 1 & 2 & 3 \\ \hline
$\stocprocremainstate=0$   & 0.3976 & 0.5901 & 0.0123 & 0 \\ \hline
$\stocprocremainstate=1$   & 0.1880 & 0.7583 & 0.0537 & 0 \\ \hline
$\stocprocremainstate=2$   & 0.0701 & 0.7676 & 0.1618 & 0.0005 \\ \hline
$\stocprocremainstate=3$   & 0.0308 & 0.6113 & 0.3559 & 0.002 \\ \hline
\textbf{$P(\roundbatches=\drawnrvdemand)$} & \textbf{0.171625} &  \textbf{0.681825} & \textbf{0.145925} &  \textbf{0.000625} \\ \hline
\end{tabular}
}
\end{table}

The distributions in Table \ref{tab:dembatches} follow directly from Eq. \eqref{eq:Az} and Eq. \eqref{eq:batcharr}. For example, when $\stocprocremainstate=0$, we know that the queue comprises an integer number of batches. Therefore, $P(\roundbatches(0) = 0) = p_{\rvdemand}(0) + p_{\rvdemand}(1) + p_{\rvdemand}(2) + p_{\rvdemand}(3)$ (Table \ref{tab:demprob}) as up to 3 new arrivals will not suffice to generate a new batch of $Q=4$ patients. Analogously, $\roundbatches(0)=1$ if we observe from 5 to 7 arrivals and $\roundbatches(0)=2$ when 8 to 10 patients arrive in the period, as these will suffice for 2 new batches. 

Conversely, when $\stocprocremainstate=3$, we have three patients in excess of a discrete number of batches, which means that $P(\roundbatches(3) = 0) = p_{\rvdemand}(0)$ as new batches will be formed whenever one or more patients arrive. One new batch will form if $\rvdemand \in [1,4]$, two new batches if $\rvdemand \in [5,8]$ and three new batches when $\rvdemand \in [8,10]$. Specifically, when 10 patients arrive, they will form a set of 13 patients when joined to the $\stocprocremainstate=3$ remaining patients from the start of the period. These 3 patients will then suffice to form three batches of $\batchsizeq=4$ patients, with one remaining batchless for the following period.

Note that $\mathbb{E}(\roundbatches) = 0.9755$, and that this value is equal to $\rho = \dfrac{\mathbb{E}(\rvdemand)}{\batchsizeq} = \dfrac{3.9022}{4} = 0.9755$. This is expected, as the dynamics of process $\stocprocbatches{\stocprocperiod}$ describes the original system, only looking at it in terms of batches of patients. Note that $\mathbb{E}(\roundbatches)$ is simply the expected number of batches of $\batchsizeq$ patients arriving and since we process $\mu=1$ batch per time period, we have $\rho = \mathbb{E}(\roundbatches)$ when looking at the evolution of the batch process $\stocprocbatches{\stocprocperiod}, \, \stocprocperiod \ge 0$.

Now, let us suppose that $\waittimereq = 6$ and $\conflevel = 0.05$ in~\eqref{eq:waitthreshold}, which means that the system requires 95\% of the patients to be operated in up to six months. To verify whether that is satisfied, we need firstly to solve Eq. \eqref{eq:ssY} for the steady state distribution of $\stocprocbatches{\stocprocperiod}, \, \stocprocperiod \ge 0$ and then use it to verify~\eqref{eq:waitthresholdy}. From Table~\ref{tab:dembatches}, we obtain $k_0=0.171625$,  $k_1=0.681825$, $k_2 = 0.145925$ and $k_3=0.000625$, which can be applied in Eq.~\eqref{eq:ssY} to find $\steadystate{\stocprocbatches{}}(0) = 0.02445$, $\steadystate{\stocprocbatches{}}(1) = 0.1182$, $\steadystate{\stocprocbatches{}}(2) = 0.1219$, $\steadystate{\stocprocbatches{}}(3) = 0.1046$, $\steadystate{\stocprocbatches{}}(4) = 0.0898$, $\steadystate{\stocprocbatches{}}(5) = 0.077$, $\steadystate{\stocprocbatches{}}(6) = 0.0661$. This yields:
\[
P(\omega \le 6 | \stocprocbatches{\stocprocperiod} \ge 1) \approx 0.602 < 0.95 = 1 - \alpha.
\]
%
This implies that $\underline Q > 4$ in Eq. \eqref{eq:under} and we should, therefore, recalculate the probability above for increasing values of $Q$ until we find the first value that satisfies the waiting time criterion - which will give us the lower bound $\underline Q$.

To verify the waiting time requirements when $\stocprocbatches{\stocprocperiod} = 0$, one can apply \eqref{eq:upperbound}, which implies that $P(\omega \le 6 | \stocprocbatches{\stocprocperiod} = 0) = 0.171625^6 \approx 2.5 \cdot 10^{-5} < 0.05$. Therefore, to find the bound $\overline Q$ in Eq. \eqref{eq:overbound}.

When verifying Eq.~\eqref{eq:under} with $Q = 10$, it implies $P(\omega \le 6 | \stocprocbatches{\stocprocperiod} = 0) = 0.609780^6 \approx 0.05114 > 0.05$. Then, if one set $\batchsizeq = 9$, it implies $P(\omega \le 6 | \stocprocbatches{\stocprocperiod} = 0) = 0.566478^6 \approx 0.03304 < 0.05$. Thus, we can set $\upq = 9$.

For the $\lowq$, one can verify $\batchsizeq = 5$, which implies $k_0=0.265720$,  $k_1=0.688120$, $k_2 = 0.046160$ and $k_3=0$, which can be applied in Eq. \eqref{eq:ssY} to find $\steadystate{\stocprocbatches{}}(0) = 0.21956$, $\steadystate{\stocprocbatches{}}(1) = 0.606723$, $\steadystate{\stocprocbatches{}}(2) = 0.143539$, $\steadystate{\stocprocbatches{}}(3) = 0.024935$, $\steadystate{\stocprocbatches{}}(4) = 0.004332$, $\steadystate{\stocprocbatches{}}(5) = 0.000752$, $\steadystate{\stocprocbatches{}}(6) = 0.000131$. This yields:
\[
P(\omega \le 6 | \stocprocbatches{\stocprocperiod} \ge 1) \approx 0.999973 > 0.95 = 1 - \alpha.
\]
%
Thus, we can set $\lowq = 5$.
%
%
\end{ex}